\ifdefined\shadowtowervtwoincluded
\else
\documentclass[11pt]{article}
\usepackage[localtheorems]{raeez-math-template}

\usepackage{array}

\theoremstyle{plain}
\newtheorem{theorem}{Theorem}[section]
\newtheorem{proposition}[theorem]{Proposition}
\newtheorem{corollary}[theorem]{Corollary}
\newtheorem{lemma}[theorem]{Lemma}
\newtheorem{conjecture}[theorem]{Conjecture}
\theoremstyle{definition}
\newtheorem{definition}[theorem]{Definition}
\newtheorem{construction}[theorem]{Construction}
\newtheorem{example}[theorem]{Example}
\theoremstyle{remark}
\newtheorem{remark}[theorem]{Remark}

\providecommand{\fg}{\mathfrak{g}}
\providecommand{\fh}{\mathfrak{h}}
\providecommand{\cA}{\mathcal{A}}
\providecommand{\cH}{\mathcal{H}}
\providecommand{\cM}{\mathcal{M}}
\providecommand{\cW}{\mathcal{W}}
\providecommand{\Sh}{\mathrm{Sh}}
\providecommand{\Res}{\operatorname{Res}}
\providecommand{\Vir}{\mathrm{Vir}}

\title{The G/L/C/M classification and shadow tables}
\author{Raeez Lorgat}
\date{}

\begin{document}
\maketitle
\fi
\providecommand{\Walg}{\mathcal{W}}

%% Four shadow classes and shadow tables

\begin{remark}[Three independent depths]
The inequalities $p_{\max} \neq k_{\max} \neq r_{\max}$ separate
OPE pole order, collision depth, and shadow depth. The shadow
classification is the stratification by the third invariant; the
operator-order trichotomy is the stratification by the first two.
\end{remark}

%% =====================================================================
%% §5. The Four Shadow Classes
%% =====================================================================

\section{The Four Shadow Classes}\label{sec:four-classes}

The single-line Riccati dichotomy partitions noncritical
one-dimensional primary lines into the three rank-one outcomes
G, L, and M. The fourth standard shadow class, C, is produced by
charged-stratum separation in the $bc$ and $\beta\gamma$ families.
The rank-one outcomes are determined by two invariants of the shadow metric
$Q_L(t) = (2\kappa + 3\alpha t)^2 + 2\Delta\, t^2$:
the cubic coefficient~$\alpha = S_3$ and the critical
discriminant $\Delta = 8\kappa S_4$.

\begin{definition}[Shadow classes]\label{def:shadow-classes}
Let $\cA$ be a chirally Koszul algebra with shadow metric $Q_L$ on a
one-dimensional primary line $L$. The \emph{shadow class} of $\cA$
along $L$ is:
\[
\renewcommand{\arraystretch}{1.4}
\begin{array}{c|cc|cc}
\textbf{Class} & \alpha & \Delta & r_{\max} & mechanism \\
\hline
G\;\text{(Gaussian)} & 0 & 0 & 2 & constant metric \\
L\;\text{(Lie/tree)} & \neq 0 & 0 & 3 & square metric \\
C\;\text{(Contact)} & 0 & \text{line data} & 4 & charged contact stratum \\
M\;\text{(Mixed)} & \neq 0 & \neq 0 & \infty & irreducible quadratic
\end{array}
\]
Here $r_{\max}$ is the \emph{shadow depth}: the largest degree~$r$
for which the shadow coefficient $S_r \neq 0$.
Class~C is the finite contact stratum of the standard $bc_\lambda$ and
$\beta\gamma_\lambda$ families, where $S_3=0$, $S_4=-5/12$, and
$S_r=0$ for $r\geq5$ on the standard conformal-weight channel.
\end{definition}

The classification is exhaustive: every entry in Table~\ref{tab:master-shadow}
falls into exactly one class, and every class is realized by at least
one standard family. The mechanism is algebraic: the shadow generating
function $H(t) = t^2\sqrt{Q_L(t)}$ is algebraic of degree~$2$
on each rank-one primary line, and its Taylor coefficients $S_r$
are computed from the convolution recursion
$H(t)^2 = t^4 Q_L(t)$, where $H(t) = \sum_{r \geq 2} r\, S_r\, t^r$
and $S_r = a_{r-2}/r$ with $a_n$ the Taylor coefficients of $\sqrt{Q_L}$.
When $Q_L$ is a perfect square ($\Delta = 0$), the square root
terminates; when $Q_L$ is irreducible ($\Delta \neq 0$), the Taylor
expansion is infinite.

\medskip

\subsection{Class G: Gaussian ($r_{\max} = 2$)}

\begin{definition}\label{def:class-G}
A chirally Koszul algebra $\cA$ is \emph{Gaussian} if
$S_r(\cA) = 0$ for all $r \geq 3$. Equivalently,
$\alpha = 0$ and $\Delta = 0$, so the shadow metric
$Q_L(t) = 4\kappa^2$ is constant and the shadow obstruction tower
terminates at the scalar curvature.
\end{definition}

\noindent\textbf{Examples.} Heisenberg $\cH_k$ with $\kappa = k$;
lattice VOAs $V_\Lambda$ with $\kappa = \operatorname{rank}(\Lambda)$;
free fermion $\psi$ with $\kappa = 1/4$.

\noindent\textbf{Characterization.}
Class~G algebras are freely generated as chiral algebras. Their bar
complex is quadratic (the differential has no higher-order terms),
and the OPE on each primary line is abelian. The full obstruction
is captured by a single scalar:
\[
\Theta_\cA = \kappa(\cA) \cdot \eta \otimes \Lambda,
\]
where $\eta$ is the cyclic direction and $\Lambda \in H^*(\bar{\cM}_g)$
is the tautological class. All genus-$g$ data is controlled by
$F_g(\cA) = \kappa(\cA) \cdot \lambda_g^{\mathrm{FP}}$ (uniform-weight).

\medskip

\subsection{Class L: Lie/tree ($r_{\max} = 3$)}

\begin{definition}\label{def:class-L}
A chirally Koszul algebra $\cA$ is \emph{Lie/tree} if
$\alpha \neq 0$ and $\Delta = 0$, so that $S_3 \neq 0$
but $S_r = 0$ for all $r \geq 4$. The shadow metric is
a perfect square:
$Q_L(t) = (2\kappa + 3\alpha t)^2$.
\end{definition}

\noindent\textbf{Examples.} Affine Kac--Moody $V_k(\fg)$ for every
simple Lie algebra $\fg$ and every non-critical level
$k \neq -h^\vee$: type~$A$ ($\mathfrak{sl}_N$),
types $B$, $C$, $D$ (orthogonal and symplectic),
exceptional types $E_6$, $E_7$, $E_8$, $G_2$, $F_4$.
For all of these:
\[
\kappa\bigl(V_k(\fg)\bigr)
= \frac{\dim(\fg)\,(k + h^\vee)}{2h^\vee}\,.
\]

\noindent\textbf{Characterization.}
The nonzero cubic shadow $S_3 = \alpha$ arises from the Lie bracket
(structure constants $f_{abc}$), while the Jacobi identity forces the
quartic shadow $S_4 = 0$ on the primary line. In trace form,
\[
r^{\mathrm{KM}}_{\mathrm{tr}}(z)=k\,\Omega_{\mathrm{tr}}/z .
\]
The KZ-normalized connection uses
\[
r^{\mathrm{KZ}}(z)
=\Omega_{\mathrm{KZ}}/((k+h^\vee)z),
\qquad
\Omega_{\mathrm{KZ}}=2h^\vee\Omega_{\mathrm{tr}}.
\]
These are conventionally related normalizations with different
$k$-dependence. The tree-level graph sum terminates at trivalent
vertices.

\medskip

\subsection{Class C: Contact ($r_{\max} = 4$)}

\begin{definition}\label{def:class-C}
A chirally Koszul algebra $\cA$ is \emph{contact} if
$S_3=0$, $S_4\neq0$, and $S_r=0$ for $r\geq5$ on the standard
contact channel. In the verified standard families this is the
conformal-weight channel of $bc_\lambda$ and
$\beta\gamma_\lambda$, where $S_4=-5/12$.
\end{definition}

\noindent\textbf{Mechanism.}
Class~C is a finite contact phenomenon external to the rank-one
Riccati dichotomy. The weight-changing line is abelian, while the
standard conformal-weight channel carries the quartic contact
coefficient. Rank-one rigidity kills the next charged target, so the
tower stops at degree~$4$.

\noindent\textbf{Examples.}
The $\beta\gamma$ system at conformal weight~$\lambda$, with
\[
\kappa(\beta\gamma_\lambda)
= 6\lambda^2 - 6\lambda + 1
= \tfrac{1}{2}\,c(\beta\gamma_\lambda),
\]
and its Koszul dual, the $bc$ system at the same weight.
At the standard weight $\lambda = 0$ or $\lambda = 1$:
$c = 2$ and $\kappa = 1$.

\medskip

\subsection{Class FF: Feigin--Frenkel locus ($\kappa = 0$)}
\label{ssec:class-FF}

The four classes~$G, L, C, M$ stratify the non-critical locus
$\kappa(\cA) \neq 0$. A fifth class appears at the boundary
$\kappa = 0$: the \emph{Feigin--Frenkel locus} where the Sugawara
topologisation fails.

\begin{definition}[Class FF]\label{def:class-FF}
A chirally Koszul algebra $\cA$ belongs to class~$\mathrm{FF}$ if
it lies at a critical Feigin--Frenkel centre or coset locus where
the Sugawara topologisation degenerates. The vanishing
$\kappa(\cA)=0$ is a necessary scalar signal in the affine examples;
the class itself is controlled by the Feigin--Frenkel coset data
$\mathrm{coset}_{FF}(\cA)$.
\end{definition}

\noindent\textbf{Examples.}
The critical-level affine algebra $V_{-h^\vee}(\fg)$ for every
simple~$\fg$ (so $V_{-2}(\mathfrak{sl}_2)$, $V_{-3}(\mathfrak{sl}_3)$,
\dots); its commutant lies in the Feigin--Frenkel centre
$\mathfrak{z}(\widehat{\fg})$, polynomial in infinitely many
Sugawara-like generators. The Koszul dual of class~L under the
involution $k \mapsto -k - 2h^\vee$ passes through class~FF at the
fixed point $k = -h^\vee$. Critical $\Walg_N$ (at the Feigin--Frenkel
image of the affine critical level) is also class~FF.

\noindent\textbf{Characterization.}
At the Feigin--Frenkel critical locus the conformal vector degenerates:
the Sugawara
construction $T = (k + h^\vee)^{-1} :J^a J^a:$ has a pole, and
topologisation~$\mathrm{SC}^{\mathrm{ch,top}} \to E_3^{\mathrm{top}}$
fails. The coset slot of the canonical fingerprint
switches from~$(\mathrm{Com}(\cH_\kappa, \cA), \cA/\mathrm{Com})$
to the Feigin--Frenkel coset
$\mathrm{coset}_{FF}(\cA) \subset \mathfrak{z}(\widehat{\fg})$.
Class~FF is the degeneration locus for the ambient shadow invariants.
On this locus the controlling datum is the Feigin--Frenkel coset data
instead of~$(\alpha, \Delta)$.

\begin{remark}[FF as a fifth companion class]
\label{rem:FF-not-refinement}
The four classes $G, L, C, M$ are determined by the shadow metric
$Q_L(t) = (2\kappa + 3\alpha t)^2 + 2\Delta t^2$, which degenerates
at $\kappa = 0$: the leading constant term vanishes, the Riccati
algebraicity argument loses its pivot, and the tower
of~$\Theta_\cA^{\leq r}$ must be read from a different source.
Together with $G,L,C,M$, this companion locus gives the
five-class fingerprint stratification used in the standard landscape.
\end{remark}

\subsection{Class M: Mixed ($r_{\max} = \infty$)}

\begin{definition}\label{def:class-M}
A chirally Koszul algebra $\cA$ is \emph{mixed} if
$\alpha \neq 0$ and $\Delta \neq 0$, so the shadow metric
$Q_L$ is an irreducible quadratic in~$t$ and the shadow obstruction tower
does not terminate.
\end{definition}

\noindent\textbf{Examples.}
Virasoro $\mathrm{Vir}_c$ with $\kappa = c/2$;
$\cW$-algebras $\cW_N$ with $\kappa = (H_N - 1)\,c$ where
$H_N = \sum_{j=1}^{N} 1/j$;
Drinfeld--Sokolov reductions of affine Kac--Moody algebras.

\noindent\textbf{Characterization.}
The shadow growth rate
\[
\rho(\cA) = \frac{\sqrt{9\alpha^2 + 16\kappa S_4}}{2|\kappa|}
\]
controls the asymptotic behaviour of $S_r$: the coefficients satisfy
$S_r \sim C(\cA)\, \rho(\cA)^r\, r^{-5/2}\, \cos(r\theta + \phi)$
for a branch-point argument~$\theta$ determined by the shadow metric.
The critical central charge $c^*$ where $\rho(\mathrm{Vir}_c) = 1$ is the
unique positive real root of the cubic
\begin{equation}\label{eq:critical-cubic}
5c^3 + 22c^2 - 180c - 872 = 0,
\qquad
c^* \approx 6.125\,.
\end{equation}
For $c > c^*$ the tower converges ($\rho < 1$), and for $c < c^*$ it
diverges ($\rho > 1$). At the self-dual point $c = 13$:
$\rho(\mathrm{Vir}_{13}) \approx 0.467$.

The Virasoro shadow obstruction tower has explicitly computed coefficients through
degree~$6$:
\begin{align}
S_2 &= \tfrac{c}{2}\,, &
S_3 &= 2\,, &
S_4 &= \frac{10}{c(5c+22)}\,,
\label{eq:vir-S234} \\
S_5 &= \frac{-48}{c^2(5c+22)}\,, &
S_6 &= \frac{4(240c+1031)}{c^3(5c+22)^2}\,.
\label{eq:vir-S567}
\end{align}
\noindent
The identity $S_3 = 2$ is independent of $c$ by direct residue
extraction from the Virasoro OPE.

\medskip

\begin{remark}[Depth and Koszulness]\label{rem:depth-vs-koszulness}
Shadow depth $r_{\max}$ does \emph{not} characterize chiral Koszulness.
All four classes (Gaussian, $r_{\max} = 2$; Lie, $r_{\max} = 3$;
contact, $r_{\max} = 4$; mixed, $r_{\max} = \infty$) consist
of chirally Koszul algebras. Shadow depth classifies \emph{complexity}
within the Koszul world: the operadic complexity conjecture identifies
$r_{\max}(\cA)$ with the $A_\infty$-depth and $L_\infty$-formality level.
\end{remark}

\begin{remark}[Total depth decomposition]\label{rem:total-depth}
The total shadow depth decomposes as
$d(\cA) = 1 + d_{\mathrm{arith}} + d_{\mathrm{alg}}$, where the algebraic depth
$d_{\mathrm{alg}} \in \{0, 1, 2, \infty\}$ and depths $\geq 5$ are
purely arithmetic (cusp forms contribute beyond the algebraic engine).
\end{remark}

\subsection{The quadrichotomy with logarithmic and Wakimoto
companions}
\label{ssec:quadrichotomy-with-companions}

The four classes $G, L, C, M$ together with the Feigin--Frenkel
locus~FF form the fingerprint stratification of the non-logarithmic
standard landscape. Two further labels occur on enlarged surfaces:
logarithmic $\cW(p)$ behaviour and Wakimoto free-field realisations.
They are companion overlays that may overlap the five strata.

\begin{definition}[Class~$M^\ast$: logarithmic companion]
\label{def:class-M-star}
A chirally Koszul algebra $\cA$ carries the companion label
$M^\ast$ if it is logarithmic in the sense of
Gurarie--Flohr~\cite{Gurarie1993,Flohr1996}, i.e.\ the Zhu
algebra $A(\cA)$ is finite-dimensional but the Hochschild Massey
products $\langle \Omega, \Omega, \ldots, \Omega \rangle$ carry
amplitude on arbitrarily high tensor powers of the vacuum.
\end{definition}

\noindent\textbf{Examples.} Triplet algebras $\cW(p)$ for $p \geq
2$; the symplectic-fermion pair at central charge $c = -2$;
logarithmic parafermion cosets.

\noindent\textbf{Characterization.}
Class~$M^\ast$ algebras satisfy $\kappa \neq 0$ and share the
mixed-shadow expansion $\alpha \neq 0$, $\Delta \neq 0$ of
class~M, but are distinguished by logarithmic amplitudes:
Zhu-finite yet Massey-unbounded. The non-logarithmic shadow
analysis applies on the formal-weight line; logarithmic corrections leave
the algebraic invariants $\alpha = S_3$ and $\Delta = 8\kappa
S_4$, which remain finite and computable, but they preclude
tempering of the shadow tower.

\begin{definition}[Class~$W$: Wakimoto stratum]
\label{def:class-W}
A chirally Koszul algebra $\cA$ belongs to class~$W$ if it
arises as the screening-kernel of a rank-$(N-1)$ Heisenberg
through Feigin--Frenkel--Wakimoto realisation,
$\cA = \bigcap_i \ker Q_{\alpha_i} \subset \cH^{\oplus(N-1)}$.
\end{definition}

\noindent\textbf{Examples.} Principal~$\cW_N$ as a Wakimoto
kernel inside $\cH^{\oplus (N-1)}$
\cite{Frenkel-Kac-Wakimoto92,FP03}; affine
$\widehat{\mathfrak{sl}}_2$ through the $\beta\gamma\,\phi$
Wakimoto realisation; admissible affine algebras through parabolic
screenings \cite{KRW03,Arakawa07}.

\noindent\textbf{Characterization.}
Class~$W$ can overlap $G,L,C,M$: it is a presentational overlay.
The ambient Heisenberg is class~$G$, while the screening
kernel may land in class~L or class~M after the BRST and
screening-obstruction data are imposed. The label records the
free-field presentation; Koszulity and Drinfeld coproduct descent
require separate screening-kernel input.

\begin{remark}[Fingerprint strata and overlays]
\label{rem:full-stratification}
The four shadow-depth classes $G, L, C, M$ stratify the
non-logarithmic non-critical locus by the algebraic invariants
$(\alpha, \Delta)$ of the shadow metric. Class~FF marks the
$\kappa = 0$ critical boundary. Class~$M^\ast$ records logarithmic
amplitude behaviour orthogonal to $(\alpha, \Delta)$: a
Zhu-finite class-M algebra is class-M, class-$M^\ast$, or both,
depending on Massey boundedness
\cite{Gurarie1993,Flohr1996}.
Class~$W$ records a presentation overlay: a class-L or class-M
algebra may simultaneously be class-$W$ when realised as a
screening kernel. The stratification is five-class:
$G,L,C,M,FF$; $M^\ast$ and $W$ are companion labels.
\end{remark}

\subsection{Leading T-line exponential base}
\label{ssec:universal-C-A-6}

The leading large-central-charge coefficients on the Virasoro
stress-tensor line have a sharp exponential base. This is a
T-line statement, separate from exact finite-$c$ radius questions for
the full shadow series of an arbitrary class-M algebra.

\begin{proposition}[Leading T-line exponential base $C^T_\cA = 6$]
\label{thm:universal-C-A-6-sa}
Let $\cA$ be a non-logarithmic class-M standard family with a
generic Virasoro stress-tensor line, and suppose the T-line
dominance hypotheses hold. Let $A_r^T$ denote the leading
large-$c$ coefficient of the T-line shadow coefficient $S_r^T$.
Then
\[
C^T_\cA \;:=\; \limsup_{r \to \infty}
\bigl|A_{r+1}^T/A_r^T\bigr| \;=\; 6.
\]
For the Virasoro line the leading generating function has the
logarithmic singularity
\[
\Sigma_{\Vir}^{T,\mathrm{lead}}(z)
\;=\;
\tfrac{4z^3}{9} - \tfrac{z^2}{9} + \tfrac{z}{27}
 - \tfrac{\log(1 + 6z)}{162}\,.
\]
\end{proposition}

\begin{proof}
The T-line large-$c$ recursion is the Virasoro Riccati recursion
with seed $S_3=2$. Solving the recursion gives the displayed
logarithmic term, whose nearest singularity is at $z=-1/6$.
The root test gives the exponential base~$6$ for the leading
coefficients $A_r^T$.
\end{proof}

For $\cW_N$ the corresponding Vol~II tempering exponent is a
separate invariant. Under the harmonic tempering law it is
\[
\beta_{\cW_N}=12(H_N-1),
\]
so $\beta_{\cW_3}=10$ and $\beta_{\cW_4}=13$. It is not the
value given by the polynomial interpolant $(N+1)(N+2)/2$.

\subsection{Picard--Fuchs spectral-curve presentation}
\label{ssec:picard-fuchs-spectral-curve}

The shadow metric $Q_L(t) = (2\kappa + 3\alpha t)^2 + 2\Delta
t^2$ is a rank-one spectral datum; its square root $\sqrt{Q_L}$
is the Seiberg--Witten-type differential of a spectral cover. On
non-logarithmic class~M, the shadow generating function
$\Sigma_\cA(z)$ satisfies a Picard--Fuchs equation whose
monodromy encodes the Borel geometry of the shadow tower.

\begin{proposition}[Picard--Fuchs equation for the leading T-line]
\label{prop:picard-fuchs-sa}
Let $\cA$ be class~M with a generic Virasoro T-line and suppose the
leading T-line dominance hypotheses of
Proposition~\ref{thm:universal-C-A-6-sa} hold. The leading T-line
shadow series $\Sigma_\cA^{T,\mathrm{lead}}(z)$ satisfies a
second-order Picard--Fuchs-type equation
\[
(1 + 6z)^2 (\Sigma_\cA^{T,\mathrm{lead}})''
 + P_1(z)\, (\Sigma_\cA^{T,\mathrm{lead}})'
 + P_0(z)\, \Sigma_\cA^{T,\mathrm{lead}}
\;=\; R_\cA(z),
\]
with $P_0, P_1, R_\cA$ polynomials of degree at most~$2$ in~$z$
determined by the leading T-line Riccati data. The regular-singular
point $z=-1/6$ is the leading T-line branch point.
\end{proposition}

\begin{proof}
Differentiate the logarithmic expression in
Proposition~\ref{thm:universal-C-A-6-sa}. Clearing the denominator
$(1+6z)^2$ gives the displayed inhomogeneous second-order equation
with polynomial coefficients. The only singularity introduced by the
logarithmic term in the leading T-line expression is $z=-1/6$.
\end{proof}

%% =====================================================================
%% §6. Shadow Tables
%% =====================================================================

\section{Shadow Tables}\label{sec:shadow-tables}

The modular characteristic $\kappa(\cA)$ is the degree-$2$ projection
of the universal MC element $\Theta_\cA$, and the genus-$g$ free energy
is $F_g(\cA) = \kappa(\cA) \cdot \lambda_g^{\mathrm{FP}}$ for
single-generator (uniform-weight) algebras at all genera, and
unconditionally at genus~$1$ for all families. Here
$\lambda_g^{\mathrm{FP}} = \frac{2^{2g-1}-1}{2^{2g-1}}
\cdot \frac{|B_{2g}|}{(2g)!}$
are the Faber--Pandharipande intersection numbers.

\subsection{Master shadow table}\label{ssec:master-shadow-table}

Table~\ref{tab:master-shadow} records, for each standard family,
the modular characteristic $\kappa$, the shadow class, the shadow
depth $r_{\max}$, the cubic coefficient $\alpha = S_3$, the quartic
coefficient $S_4$, and the critical discriminant $\Delta = 8\kappa S_4$.

\begin{table}[ht]
\centering
\caption{Master shadow table: standard landscape}
\label{tab:master-shadow}
\renewcommand{\arraystretch}{1.35}
{\small
\begin{tabular}{lcccccc}
\toprule
\textbf{Algebra $\cA$}
 & $\boldsymbol{\kappa(\cA)}$
 & \textbf{Class}
 & $\boldsymbol{r_{\max}}$
 & $\boldsymbol{S_3}$
 & $\boldsymbol{S_4}$
 & $\boldsymbol{\Delta}$ \\
\midrule
\multicolumn{7}{l}{\textit{Class G (Gaussian)}} \\[2pt]
Heisenberg $\cH_k$
 & $k$
 & G & $2$ & $0$ & $0$ & $0$ \\
Lattice $V_\Lambda$
 & $\operatorname{rank}(\Lambda)$
 & G & $2$ & $0$ & $0$ & $0$ \\
Free fermion $\psi$
 & $1/4$
 & G & $2$ & $0$ & $0$ & $0$ \\[3pt]
\midrule
\multicolumn{7}{l}{\textit{Class L (Lie/tree)}} \\[2pt]
$\widehat{\mathfrak{sl}}_2$ at level $k$
 & $3(k{+}2)/4$
 & L & $3$
 & $\neq 0$ & $0$ & $0$ \\
$\widehat{\mathfrak{sl}}_N$ at level $k$
 & $(N^2{-}1)(k{+}N)/(2N)$
 & L & $3$
 & $\neq 0$ & $0$ & $0$ \\
$\widehat{E}_6$ at level $k$
 & $13(k{+}12)/4$
 & L & $3$
 & $\neq 0$ & $0$ & $0$ \\
$\widehat{E}_7$ at level $k$
 & $133(k{+}18)/36$
 & L & $3$
 & $\neq 0$ & $0$ & $0$ \\
$\widehat{E}_8$ at level $k$
 & $62(k{+}30)/15$
 & L & $3$
 & $\neq 0$ & $0$ & $0$ \\[3pt]
\midrule
\multicolumn{7}{l}{\textit{Class C (Contact)}} \\[2pt]
$bc$ ghosts (weight $\lambda$)
 & $-(6\lambda^2{-}6\lambda{+}1)$
 & C & $4$
 & $0$ & $-\dfrac{5}{12}$
 & $\dfrac{10}{3}(6\lambda^2{-}6\lambda{+}1)$ \\
$\beta\gamma$ (weight $\lambda$)
 & $6\lambda^2{-}6\lambda{+}1$
 & C & $4$
 & $0$ & $-\dfrac{5}{12}$
 & $-\dfrac{10}{3}(6\lambda^2{-}6\lambda{+}1)$ \\[3pt]
\midrule
\multicolumn{7}{l}{\textit{Class M (Mixed)}} \\[2pt]
$\mathrm{Vir}_c$
 & $c/2$
 & M & $\infty$
 & $2$
 & $\dfrac{10}{c(5c{+}22)}$
 & $\dfrac{40}{5c{+}22}$ \\[8pt]
$\cW_3$ at $c$
 & $5c/6$
 & M & $\infty$
 & $\neq 0$ & $\neq 0$ & $\neq 0$ \\
$\cW_4$ at $c$
 & $13c/12$
 & M & $\infty$
 & $\neq 0$ & $\neq 0$ & $\neq 0$ \\
$\cW_5$ at $c$
 & $77c/60$
 & M & $\infty$
 & $\neq 0$ & $\neq 0$ & $\neq 0$ \\
$\cW_N$ at $c$ (general)
 & $(H_N{-}1)\,c$
 & M & $\infty$
 & $\neq 0$ & $\neq 0$ & $\neq 0$ \\
\bottomrule
\end{tabular}
}
\end{table}

\noindent
For $bc_\lambda$ and $\beta\gamma_\lambda$, the class-C witness is
the standard conformal-weight family:
$S_2=6\lambda^2-6\lambda+1$, $S_3=0$, $S_4=-5/12$, and
$S_r=0$ for $r\geq5$, with the sign of $\Delta=8\kappa S_4$
set by the sign of~$\kappa$.

\noindent
In the table, $H_N = 1 + 1/2 + \dotsb + 1/N$ denotes the
$N$-th harmonic number. The Virasoro quartic contact
$S_4 = Q^{\mathrm{contact}}_{\mathrm{Vir}} = 10/[c(5c+22)]$
is the unique quartic shadow on the $T$-line, proved by direct
computation from the Virasoro OPE. The T-line shadow data ($S_3 = 2$,
$S_4 = 10/[c(5c+22)]$) is \emph{universal}: it coincides for $\cW_N$
at all $N \geq 2$, since every $\cW_N$ contains the Virasoro
subalgebra along the stress-tensor primary line. The additional
primary lines (spin-$3$, spin-$4$, etc.) carry independent shadow
data with separate quartic coefficients.


\subsection{Exceptional and non-simply-laced types}\label{ssec:exceptional-table}

Table~\ref{tab:exceptional-shadow} records the affine Kac--Moody
data for exceptional and non-simply-laced types. All are class~L.

\begin{table}[ht]
\centering
\caption{Shadow data for exceptional and non-simply-laced types}
\label{tab:exceptional-shadow}
\renewcommand{\arraystretch}{1.35}
{\small
\begin{tabular}{lcccccc}
\toprule
$\fg$
 & $\dim(\fg)$
 & $h^\vee$
 & $\boldsymbol{\kappa\bigl(V_k(\fg)\bigr)}$
 & $c(V_k) + c(V_{k'})$
 & \textbf{Class}
 & $r_{\max}$ \\
\midrule
$B_2 \cong \mathfrak{so}_5$
 & $10$ & $3$ & $5(k{+}3)/3$ & $20$ & L & $3$ \\
$C_2 \cong \mathfrak{sp}_4$
 & $10$ & $3$ & $5(k{+}3)/3$ & $20$ & L & $3$ \\
$G_2$
 & $14$ & $4$ & $7(k{+}4)/4$ & $28$ & L & $3$ \\
$F_4$
 & $52$ & $9$ & $26(k{+}9)/9$ & $104$ & L & $3$ \\
$E_6$
 & $78$ & $12$ & $13(k{+}12)/4$ & $156$ & L & $3$ \\
$E_7$
 & $133$ & $18$ & $133(k{+}18)/36$ & $266$ & L & $3$ \\
$E_8$
 & $248$ & $30$ & $62(k{+}30)/15$ & $496$ & L & $3$ \\
\bottomrule
\end{tabular}
}
\end{table}

\noindent
The kappa formula is universal:
$\kappa(V_k(\fg)) = \dim(\fg)\,(k + h^\vee)/(2h^\vee)$.
The central charge sum $c + c' = 2\dim(\fg)$ under the
Feigin--Frenkel involution $k \mapsto -k - 2h^\vee$, and
the complementarity $\kappa + \kappa' = 0$ holds for all
affine Kac--Moody algebras. For non-simply-laced types,
the dual Coxeter number $h^\vee$ differs from the Coxeter
number $h$ ($h^\vee = 3$ for $B_2$ and $C_2$, while $h = 4$).
The shadow class is~L regardless: the Jacobi identity kills the
quartic on every primary line.

\begin{remark}[$B_2$/$C_2$ isomorphism]
The Lie algebras $B_2 = \mathfrak{so}_5$ and
$C_2 = \mathfrak{sp}_4$ are isomorphic, sharing $\dim = 10$,
$h^\vee = 3$, and $h = 4$. Their Cartan matrices are transposes
of each other, but all shadow invariants coincide.
\end{remark}


\subsection{Complementarity table}\label{ssec:complementarity-table}

\begin{table}[ht]
\centering
\caption{Complementarity data: $\kappa(\cA) + \kappa(\cA^!)$ and
$c(\cA) + c(\cA^!)$ under Koszul duality}
\label{tab:complementarity}
\renewcommand{\arraystretch}{1.35}
{\small
\begin{tabular}{llccc}
\toprule
\textbf{Family $\cA$}
 & \textbf{Koszul dual $\cA^!$}
 & $\boldsymbol{c + c'}$
 & $\boldsymbol{\kappa + \kappa'}$
 & \textbf{Self-dual point} \\
\midrule
$\cH_k$
 & $\mathrm{Sym}^{\mathrm{ch}}(V^*)$
 & ${}^{\ddagger}$ & $0$ & none \\
$V_\Lambda$ (lattice)
 & $V_{\Lambda^!}$
 & ${}^{\ddagger}$ & $0$ & $\Lambda$ self-dual \\
Free fermion
 & $\mathrm{Sym}^{\mathrm{ch}}(\gamma)$
 & $0$ & $0$ & none \\
$\beta\gamma_\lambda$
 & $bc_\lambda$
 & $0$ & $0$ & none \\[3pt]
\midrule
$V_k(\mathfrak{sl}_2)$
 & $V_{-k-4}(\mathfrak{sl}_2)$
 & $6$ & $0$ & $k = -2$ (critical) \\
$V_k(\mathfrak{sl}_N)$
 & $V_{-k-2N}(\mathfrak{sl}_N)$
 & $2(N^2{-}1)$ & $0$ & $k = -N$ (critical) \\
$V_k(\fg)$ (general)
 & $V_{-k-2h^\vee}(\fg)$
 & $2\dim(\fg)$ & $0$ & $k = -h^\vee$ (critical) \\[3pt]
\midrule
$\mathrm{Vir}_c$
 & $\mathrm{Vir}_{26-c}$
 & $26$ & $13$ & $c = 13$ \\
$\cW_3$ at $c$
 & $\cW_3$ at $100{-}c$
 & $100$ & $250/3$ & $c = 50$ \\
$\cW_4$ at $c$
 & $\cW_4$ at $246{-}c$
 & $246$ & $533/2$ & $c = 123$ \\
$\cW_5$ at $c$
 & $\cW_5$ at $488{-}c$
 & $488$ & $9394/15$ & $c = 244$ \\
$\cW_N$ at $c$ (general)
 & $\cW_N$ at $K_N{-}c$
 & $K_N$ & $(H_N{-}1)\,K_N$ & $c = K_N/2$ \\
\bottomrule
\end{tabular}
}
\end{table}

\noindent
${}^{\ddagger}$The Koszul dual $\cH_k^! = \mathrm{Sym}^{\mathrm{ch}}(V^*)$
has $\kappa(\cH_k^!) = -k = \kappa(\cH_{-k})$; its OPE structure
differs from that of $\cH_{-k}$.
The central charge sum has no single-parameter-family meaning.

\noindent
Here $K_N = 2(N{-}1)(2N^2{+}2N{+}1)$ is the Koszul conductor for
$\cW_N = \cW(\mathfrak{sl}_N)$ (the sum $c + c'$ under the
Feigin--Frenkel involution $k \mapsto -k - 2N$ on the underlying
$\mathfrak{sl}_N$). The key structural distinction: for all
affine Kac--Moody algebras, $\kappa + \kappa' = 0$ (anti-symmetry
under the Feigin--Frenkel involution); for $\cW$-algebras,
$\kappa + \kappa' = (H_N - 1) \cdot K_N \neq 0$.
The Virasoro ($N = 2$) gives $\kappa + \kappa' = \tfrac{1}{2} \cdot 26
= 13$, a nonzero constant independent of~$c$. The Koszul involution
$c \mapsto 26 - c$ has a fixed point at~$c = 13$, where
$\kappa(\mathrm{Vir}_{13}) = 13/2$ and $\mathrm{Vir}_{13}$ is self-dual
under uncurved quadratic duality.


\subsection{Genus tables}\label{ssec:genus-tables}

The Faber--Pandharipande numbers and the genus expansion
$F_g(\cA) = \kappa(\cA) \cdot \lambda_g^{\mathrm{FP}}$ (uniform-weight).

\begin{table}[ht]
\centering
\caption{Faber--Pandharipande intersection numbers $\lambda_g^{\mathrm{FP}}$}
\label{tab:faber-pandharipande}
\renewcommand{\arraystretch}{1.35}
{\small
\begin{tabular}{ccc}
\toprule
$g$ & $\lambda_g^{\mathrm{FP}}$ & Decimal \\
\midrule
$1$ & $1/24$ & $0.041\overline{6}$ \\
$2$ & $7/5760$ & $1.215 \times 10^{-3}$ \\
$3$ & $31/967680$ & $3.204 \times 10^{-5}$ \\
$4$ & $127/154828800$ & $8.201 \times 10^{-7}$ \\
$5$ & $73/3503554560$ & $2.084 \times 10^{-8}$ \\
\bottomrule
\end{tabular}
}
\end{table}

\noindent
The formula is
$\lambda_g^{\mathrm{FP}} = \frac{2^{2g-1} - 1}{2^{2g-1}}
\cdot \frac{|B_{2g}|}{(2g)!}$,
where $B_{2g}$ is the $2g$-th Bernoulli number. The generating
function is $\sum_{g \geq 1} \lambda_g^{\mathrm{FP}}\, \hbar^{2g}
= \hat{A}(i\hbar) - 1$, where $\hat{A}(x) = (x/2)/\sinh(x/2)$ is
the $\hat{A}$-genus (so $\hat{A}(ix) = (x/2)/\sin(x/2)$). The numbers $\lambda_g^{\mathrm{FP}}$ are
strictly positive, strictly decreasing, and decay like
$(2\pi)^{-2g}$.

\begin{table}[ht]
\centering
\caption{Genus expansion $F_g(\cA) = \kappa \cdot \lambda_g^{\mathrm{FP}}$ (uniform-weight)
for selected families}
\label{tab:genus-expansion}
\renewcommand{\arraystretch}{1.35}
{\small
\begin{tabular}{lccccc}
\toprule
\textbf{Algebra $\cA$}
 & $\boldsymbol{\kappa}$
 & $\boldsymbol{F_1}$
 & $\boldsymbol{F_2}$
 & $\boldsymbol{F_3}$ \\
\midrule
$\cH_1$ (Heisenberg, $k{=}1$)
 & $1$ & $1/24$ & $7/5760$ & $31/967680$ \\
$V_1(\mathfrak{sl}_2)$
 & $9/4$ & $3/32$ & $7/2560$ & $31/430080$ \\
$V_1(\mathfrak{sl}_3)$
 & $16/3$ & $2/9$ & $7/1080$ & $31/181440$ \\
$V_\Lambda$ (Leech, $\operatorname{rank}{=}24$)
 & $24$ & $1$ & $7/240$ & $31/40320$ \\
$\beta\gamma$ (standard, $\lambda{=}1$)
 & $1$ & $1/24$ & $7/5760$ & $31/967680$ \\
$\mathrm{Vir}_1$
 & $1/2$ & $1/48$ & $7/11520$ & $31/1935360$ \\
$\mathrm{Vir}_{13}$
 & $13/2$ & $13/48$ & $91/11520$ & $403/1935360$ \\
$\mathrm{Vir}_{26}$
 & $13$ & $13/24$ & $91/5760$ & $403/967680$ \\
$\cW_3$ at $c{=}6$
 & $5$ & $5/24$ & $7/1152$ & $31/193536$ \\
\bottomrule
\end{tabular}
}
\end{table}

\noindent
The entries for $V_1(\mathfrak{sl}_3)$ use
$\kappa = (8 \cdot 4)/(2 \cdot 3) = 16/3$. The Leech lattice
($\operatorname{rank} = 24$) gives $F_1 = 1$, reflecting the
Jacobi $\Delta$-function origin of the Leech theta series.

\subsection{Shadow growth rate table}\label{ssec:growth-rate-table}

\begin{table}[ht]
\centering
\caption{Shadow growth rate $\rho(\cA)$ for class~M algebras
(Virasoro at selected central charges)}
\label{tab:growth-rate}
\renewcommand{\arraystretch}{1.35}
{\small
\begin{tabular}{ccccl}
\toprule
$c$
 & $\kappa$
 & $Q^{\mathrm{contact}}$
 & $\rho$
 & \textbf{Regime} \\
\midrule
$1/2$ (Ising)
 & $1/4$ & $40/49$ & $\approx 12.53$ & divergent \\
$1$
 & $1/2$ & $10/27$ & $\approx 6.24$ & divergent \\
$4$
 & $2$ & $5/84$ & $\approx 1.54$ & divergent \\
$c^* \approx 6.125$
 & $\approx 3.06$ & $\approx 0.031$ & $1$ & critical \\
$13$ (self-dual)
 & $13/2$ & $10/1131$ & $\approx 0.467$ & convergent \\
$25$
 & $25/2$ & $2/735$ & $\approx 0.242$ & convergent \\
$26$ (bosonic string)
 & $13$ & $5/1976$ & $\approx 0.232$ & convergent \\
\bottomrule
\end{tabular}
}
\end{table}

\noindent
The shadow growth rate $\rho = \sqrt{(9\alpha^2 + 16\kappa S_4)/(4\kappa^2)}$
determines the convergence of the shadow obstruction tower. For Virasoro,
$\rho^2 = (180c + 872)/[c^2(5c+22)]$. The critical central charge
$c^* \approx 6.125$ (the positive real root of~\eqref{eq:critical-cubic})
separates the convergent regime ($c > c^*$, tower summable) from the
divergent regime ($c < c^*$, tower grows exponentially). Class~G and
class~L algebras have $\rho = 0$ (the tower terminates, so there is no
exponential growth).

\subsection{Virasoro shadow obstruction tower table}\label{ssec:virasoro-tower-table}

\begin{table}[ht]
\centering
\caption{Virasoro shadow coefficients $S_r$ for $r = 2, \dotsc, 10$}
\label{tab:virasoro-tower}
\renewcommand{\arraystretch}{1.35}
{\small
\begin{tabular}{cl}
\toprule
$r$ & $S_r(\mathrm{Vir}_c)$ \\
\midrule
$2$ & $c/2$ \\
$3$ & $2$ \\
$4$ & $10/[c(5c+22)]$ \\
$5$ & $-48/[c^2(5c+22)]$ \\
$6$ & $80(45c+193)/[3c^3(5c+22)^2]$ \\
$7$ & $-2880(15c+61)/[7c^4(5c+22)^2]$ \\
$8$ & $80(2025c^2+16470c+33314)/[c^5(5c+22)^3]$ \\
$9$ & $-1280(2025c^2+15570c+29554)/[3c^6(5c+22)^3]$ \\
$10$ & $256(91125c^3+1050975c^2+3989790c+4969967)/[c^7(5c+22)^4]$ \\
\bottomrule
\end{tabular}
}
\end{table}

\noindent
Each $S_r$ is a rational function of $c$ with poles at $c = 0$ and
$c = -22/5$; these are the two values where the quadratic Casimir
$c(5c+22)$ vanishes. The sign pattern
$S_r > 0$ for even $r$ and $S_r < 0$ for odd $r \geq 5$
(at $c > 0$) reflects the oscillatory decay $\cos(r\theta + \phi)$
from the complex branch points of the shadow generating function.
All entries have been verified by two independent methods: closed-form
derivation and the convolution recursion from $H^2 = t^4 Q_L$.


%% =====================================================================
%% §7. Structural Theorems
%% =====================================================================

\section{Structural Theorems}\label{sec:structural-theorems}

The classification is supported by four structural statements.
Proposition~\ref{prop:depth-gap-standalone} excludes intermediate
algebraic depths; Proposition~\ref{prop:sc-formality-standalone}
identifies class~$\mathbf{G}$ with SC-formality; and
Proposition~\ref{prop:betagamma-contact-standalone} produces
class~$\mathbf{C}$ by contact-channel separation. The
Drinfeld--Sokolov conjecture isolates the class-raising functor under
its generic-level hypotheses, and
Proposition~\ref{prop:shadow-semilattice-standalone}
gives the ordered classes their split direct-sum join.


\subsection{The depth gap theorem}\label{ssec:depth-gap}

The Riccati equation leaves only three rank-one possibilities. The
tower has support through $S_2$ in class~$\mathbf{G}$, through $S_3$
in class~$\mathbf{L}$, through a contact-channel $S_4$ in
class~$\mathbf{C}$, and through infinitely many nonzero coefficients
in class~$\mathbf{M}$. A quadratic $Q_L$ is either a square or a
non-square; its square root is respectively polynomial or an infinite
binomial series. Thus the discriminant of $Q_L$ determines the
rank-one shadow depth, while class~$\mathbf{C}$ is supplied by the
separate standard contact channel.

Two distinct notions of depth appear in this analysis and must
be kept separate. The algebraic depth
$d_{\mathrm{alg}}(\cA)$ counts independent shadow coefficients:
$d_{\mathrm{alg}} = 0$ if $S_r = 0$ for all $r \geq 3$,
$d_{\mathrm{alg}} = n$ if $S_r = 0$ for all $r \geq n + 3$,
and $d_{\mathrm{alg}} = \infty$ if infinitely many $S_r$ are
nonzero. The generating depth $d_{\mathrm{gen}}(\cA)$ is the
degree at which higher operations are recursively determined by
lower ones. For the Virasoro algebra, $d_{\mathrm{gen}} = 3$
(the cubic $S_3 = 2$ and quartic $S_4 = 10/[c(5c{+}22)]$
determine all $S_r$ via the convolution recursion) but
$d_{\mathrm{alg}} = \infty$ (every $S_r$ is nonzero). The
recursion is finitely generated; the output has infinite support.

\begin{proposition}[Depth gap]\label{prop:depth-gap-standalone}
Let $\cA$ be a chirally Koszul algebra in the standard landscape
with $\kappa(\cA) \neq 0$. Then
\[
 d_{\mathrm{alg}}(\cA) \;\in\; \{0,\, 1,\, 2,\, \infty\}.
\]
No finite value $d_{\mathrm{alg}} \geq 3$ is realized.
The four values correspond bijectively to the shadow classes
$\mathbf{G}$, $\mathbf{L}$, $\mathbf{C}$, $\mathbf{M}$.
\end{proposition}

\begin{proof}
The argument splits into a rank-one analysis on each primary
line and a multi-sector extension.

\emph{Rank one.}
On a primary line~$L$ with $\kappa \neq 0$, the shadow generating
function $H(t) = t^2 \sqrt{Q_L(t)}$ is determined by
$(\kappa, \alpha, S_4)$, where the shadow metric
$Q_L(t) = (2\kappa + 3\alpha\, t)^2 + 2\Delta\, t^2$
is the Gaussian decomposition and
$\Delta = 8\kappa S_4$ is the critical discriminant.

If $\Delta = 0$ (equivalently $S_4 = 0$), then
$Q_L(t) = (2\kappa + 3\alpha\, t)^2$ is a perfect square, and
$H(t) = t^2(2\kappa + 3\alpha\, t)$ is a polynomial of
degree~$3$. The shadow coefficients $S_r = 0$ for all
$r \geq 4$, giving $d_{\mathrm{alg}} \in \{0, 1\}$ according
to whether $\alpha = 0$ (class~$\mathbf{G}$) or
$\alpha \neq 0$ (class~$\mathbf{L}$).

If $\Delta \neq 0$ (equivalently $S_4 \neq 0$), then $Q_L$ is
irreducible over the rank-one coefficient field: the interaction correction
$2\Delta\, t^2$ prevents $\sqrt{Q_L}$ from being polynomial.
The Taylor expansion of $\sqrt{Q_L(t)}$ about $t = 0$ has
infinitely many nonzero coefficients (the binomial series for
$(1 + u)^{1/2}$ with
$u = (6\alpha\,t + (9\alpha^2 + 2\Delta)\,t^2)/(4\kappa^2)$
has infinitely many nonzero terms when the coefficient of~$t^2$ in~$u$
is nonzero). Therefore $S_r \neq 0$ for infinitely many~$r$,
and $d_{\mathrm{alg}} = \infty$ (class~$\mathbf{M}$).

The gap at $d_{\mathrm{alg}} = 3$ is now forced. Suppose
$S_4 \neq 0$. The convolution recursion at degree~$5$ gives
$S_5 = -(6/5)(1/\kappa)\,\alpha\,S_4$.
If $\alpha \neq 0$: this is nonzero, so $S_5 \neq 0$.
If $\alpha = 0$: then $S_5 = 0$, but the degree-$6$ recursion
gives $S_6 \propto S_4^2/\kappa \neq 0$. In either case,
$S_4 \neq 0$ propagates to infinitely many nonzero $S_r$: the
closed form $H(t) = t^2\sqrt{Q_L(t)}$ makes this manifest.
Termination at any finite degree $\geq 5$ is impossible.

Class~$\mathbf{C}$ ($d_{\mathrm{alg}} = 2$) occupies the
contact channel where $S_3=0$, $S_4=-5/12$, and $S_r=0$ for
$r\geq5$ by a mechanism external to the rank-one Riccati analysis. See
Proposition~\ref{prop:betagamma-contact-standalone} below.

\emph{Multi-sector extension.}
For a multi-generator algebra with primary lines
$L_1, \dotsc, L_n$, the cubic pump propagates: if
$\alpha \neq 0$ on any mixed sector, the adjoint
$\mathrm{ad}_{S_3}$ generates an infinite chain on every
reachable sector, forcing $d_{\mathrm{alg}} = \infty$. If all
mixed cubics vanish, the sectors decouple and
$d_{\mathrm{alg}} \leq \max_i d_{\mathrm{alg}}(L_i) \leq 2$.
No intermediate value $d_{\mathrm{alg}} = 3$ is accessible.
\end{proof}


\subsection{SC-formality characterises class G}
\label{ssec:sc-formality}

Every algebra in the standard landscape is chirally Koszul: the
bar spectral sequence collapses at $E_2$, the bar cohomology
concentrates in degree~$1$, and the transferred $A_\infty$
operations on $H^*(\bar{B}^{\mathrm{ch}}(\cA))$ satisfy
$m_k = 0$ for $k \geq 3$. Koszulness is common to all four shadow
classes.

A stronger condition does. Swiss-cheese formality asks that the
higher SC-operations on $\cA$ \emph{itself} vanish:
$m_k^{\mathrm{SC}} = 0$ for all $k \geq 3$. This is a
condition on the algebra, distinct from a condition on its bar cohomology; the common
notation $m_k$ for both is a source of systematic confusion.
The question is: which algebras satisfy this stronger condition?

The answer is a unique-survivor argument. Four shadow classes
enter; one constraint eliminates three.

\begin{proposition}[SC-formality iff class~$\mathbf{G}$]
\label{prop:sc-formality-standalone}
Let $\cA$ be a chiral algebra in the standard landscape.
Then $\cA$ is Swiss-cheese formal if and only if $\cA$ belongs
to class~$\mathbf{G}$.
\end{proposition}

\begin{proof}
The forward direction is immediate.
If $\cA$ is class~$\mathbf{G}$ (Heisenberg, lattice, free
fermion), the current algebra bracket is central:
$[a_m, a_n] = km\,\delta_{m+n,0}$ for Heisenberg, and the
analogous statement for every other class~$\mathbf{G}$ family.
All nested brackets vanish, so $m_k^{\mathrm{SC}} = 0$ for
$k \geq 3$.

Conversely, suppose $\cA$ is SC-formal. The cubic shadow $S_3$ is the normalized projection
of the trilinear form
\[
 C(x,y,z) \;=\; \kappa\bigl(x,\, [y,z]\bigr),
\]
where $\kappa$ is the invariant bilinear form and $[-,-]$ the
current algebra bracket. The chain of implications is:
\[
 \text{SC-formal}
 \;\Longrightarrow\;
 m_3^{\mathrm{SC}} = 0
 \;\Longrightarrow\;
 S_3 = 0
 \;\Longrightarrow\;
 C = 0.
\]
The first arrow is the definition. The second holds because
$S_3$ is the degree-$3$ shadow projection of $m_3^{\mathrm{SC}}$.
The third requires the key input: $\kappa$ is nondegenerate on
the standard landscape. With $\kappa$ nondegenerate, $C = 0$
forces $[y,z] = 0$ for all $y, z$. The current algebra is
abelian; an abelian current algebra at level~$k$ is Heisenberg;
Heisenberg is class~$\mathbf{G}$.

Classes~$\mathbf{L}$, $\mathbf{C}$, and~$\mathbf{M}$ all have
$S_3 \neq 0$ or carry structure (charged strata, nonlinear
composites) incompatible with $m_3^{\mathrm{SC}} = 0$. Each is
eliminated. Only class~$\mathbf{G}$ survives.
\end{proof}

\begin{remark}[The two meanings of formality]
\label{rem:sc-vs-ainfty-standalone}
The distinction between $A_\infty$-formality of bar cohomology
and SC-formality of the algebra itself is the central subtlety
of the Koszulness programme. Every chirally Koszul algebra is
$A_\infty$-formal on its bar cohomology; this is a consequence of
the $E_2$-collapse and holds in all four shadow classes.
SC-formality is a strictly stronger condition, operative on the
algebra rather than on its cohomological shadow.
Proposition~\ref{prop:sc-formality-standalone} identifies the
precise locus where the two notions coincide: class~$\mathbf{G}$.
\end{remark}


\subsection{The betagamma contact mechanism}
\label{ssec:betagamma-contact}

The single-line Riccati dichotomy produces three outcomes:
$d_{\mathrm{alg}} \in \{0, 1, \infty\}$. Class~$\mathbf{C}$
($d_{\mathrm{alg}} = 2$) is absent from this list. Its
existence requires a mechanism that the rank-one analysis
misses: shadow data that lives on a charged stratum rather
than on the primary line, and a termination mechanism that
operates by dimensional obstruction rather than by
algebraic identity.

The $\beta\gamma$ system is the algebra that exhibits this
mechanism. It is the unique standard family where shadow depth
exceeds collision depth: the OPE pole order is $p_{\max} = 1$
(a single simple pole $\beta(z)\gamma(w) \sim 1/(z{-}w)$), the
collision depth is $k_{\max} = 0$ (no collision singularity),
yet the shadow depth is $r_{\max} = 4$. Three invariants that
agree for every other standard family here separate:
$p_{\max} \neq k_{\max} \neq r_{\max}$. The mechanism is
stratum separation: the quartic shadow lives on a charged
stratum invisible to the single-line dichotomy, and the tower
terminates there by rank-one rigidity.

\begin{proposition}[Contact class via stratum separation]
\label{prop:betagamma-contact-standalone}
Let $\beta\gamma_\lambda$ denote the $\beta\gamma$ system at
conformal weight~$\lambda$, with central charge
$c(\beta\gamma_\lambda) = 2(6\lambda^2 - 6\lambda + 1)$ and
$\kappa(\beta\gamma_\lambda) = 6\lambda^2 - 6\lambda + 1$.
Then\textup{:}
\begin{enumerate}[label=\textup{(\roman*)}]
\item On the weight-changing primary line, the OPE is abelian
 \textup{(}$\alpha = 0$\textup{)} and $\Delta = 0$:
 the single-line dichotomy gives $r_{\max}|_L = 2$
 \textup{(}class~$\mathbf{G}$ on this line\textup{)}.
\item The quartic contact invariant
 $\mathfrak{Q}_{\beta\gamma} = \mathrm{cyc}(m_3)$ is nonzero
 and lives in the charged sector
 $\mathrm{Def}_{\mathrm{cyc}}^{(q)}$
 with $q \neq 0$.
\item The self-bracket
 $\mathrm{Def}_{\mathrm{cyc}}^{(q)} \otimes
 \mathrm{Def}_{\mathrm{cyc}}^{(q)} \to
 \mathrm{Def}_{\mathrm{cyc}}^{(2q)}$
 has $\dim \mathrm{Def}_{\mathrm{cyc}}^{(2q)} = 0$ for
 rank-one systems: the quartic pump has zero target, and
 the tower terminates at $r_{\max} = 4$.
\item The shadow depth is $r_{\max}(\beta\gamma_\lambda) = 4$
 and $d_{\mathrm{alg}}(\beta\gamma_\lambda) = 2$ for all
 $\lambda$.
\end{enumerate}
\end{proposition}

\begin{proof}
The primary-line computation gives (i), the charged-sector
decomposition gives (ii), rank-one rigidity gives (iii), and
(i)--(iii) imply (iv).

\emph{Item~(i): invisibility on the primary line.}
The $\beta\gamma$ OPE
$\beta(z)\gamma(w) \sim 1/(z{-}w)$ is a simple pole.
The self-OPEs $\beta(z)\beta(w)$ and $\gamma(z)\gamma(w)$
are regular. There is no current algebra bracket: the OPE on
the weight-changing line is abelian, $\alpha = 0$, $S_4 = 0$,
$\Delta = 0$. The Riccati dichotomy applied to this line
gives $d_{\mathrm{alg}} = 0$. If the primary line were the
whole story, $\beta\gamma$ would be class~$\mathbf{G}$.

\emph{Item~(ii): birth of the quartic on the charged stratum.}
The weight-changing primary line is incomplete. The $\beta\gamma$
system has two generators of different conformal weight: the
cyclic deformation complex decomposes by charge, and the charged
sector $\mathrm{Def}_{\mathrm{cyc}}^{(q)}$ with $q \neq 0$ is
invisible to the restriction $q = 0$ that defines the primary
line. The quartic contact invariant
$\mathfrak{Q}_{\beta\gamma} = \mathrm{cyc}(m_3)$ is the
trilinear Massey product on this charged stratum. In the standard
conformal-weight family the quartic coefficient is
$S_4=-5/12$ for all~$\lambda$; at $\lambda=1/2$ the symplectic
boson changes the parity grading, while the graded contact witness
remains nonzero.

\emph{Item~(iii): death of the quintic by rank-one rigidity.}
The quartic shadow, once born, must either propagate (producing
$d_{\mathrm{alg}} = \infty$, class~$\mathbf{M}$) or die (giving
$d_{\mathrm{alg}} = 2$, class~$\mathbf{C}$). The propagation
mechanism is the quartic pump: the self-bracket
$\mathrm{Def}_{\mathrm{cyc}}^{(q)} \otimes
\mathrm{Def}_{\mathrm{cyc}}^{(q)} \to
\mathrm{Def}_{\mathrm{cyc}}^{(2q)}$.
For rank-one systems (one generator in each charge sector), the
target space has dimension zero: there is no field of charge~$2q$
in a rank-one system. The quintic obstruction $S_5$ vanishes
identically on the charged stratum; the pump has no room to
operate. The tower terminates.

\emph{Item~(iv).}
Items~(i)--(iii) give $S_r = 0$ for $r \geq 5$ on every stratum
and $S_4 \neq 0$ on the charged stratum: $r_{\max} = 4$,
$d_{\mathrm{alg}} = 2$. The result holds for all $\lambda$ (the
rank-one rigidity is independent of conformal weight), confirming
class~$\mathbf{C}$.
\end{proof}

\begin{remark}[Why contact is rare]
\label{rem:contact-rarity}
The mechanism that produces class~$\mathbf{C}$ requires three
simultaneous conditions: an abelian primary line ($\alpha = 0$),
a nonzero quartic on a charged stratum
($\mathfrak{Q} \neq 0$), and rank-one rigidity to terminate
the tower ($\dim\mathrm{Def}^{(2q)} = 0$). The third condition
fails for rank $\geq 2$: the target of the quartic pump acquires
positive dimension, and the tower generically diverges. Among
standard families, only the $\beta\gamma$ system and its Koszul
dual $bc$ satisfy all three conditions.
\end{remark}


\subsection{Shadow depth under Drinfeld--Sokolov reduction}
\label{ssec:ds-shadow-depth}

The four shadow classes are static invariants of individual
algebras. A natural question is how these invariants
transform under the functors that connect families. Quantum
Drinfeld--Sokolov reduction is the principal such functor: it
sends an affine Kac--Moody algebra $V_k(\fg)$ to a
$\cW$-algebra $\cW^k(\fg, f)$, parameterised by a choice of
nilpotent $f \in \fg$. At $f = 0$ the reduction is trivial
and the output is the input. At $f \neq 0$ the output is a
genuinely different algebra: new composite fields appear (the
Sugawara stress tensor $T$, higher-spin currents $W_s$), and
these composites carry OPE data absent from the affine algebra.

The conjectural shadow-theoretic consequence is class escalation. An
affine Kac--Moody algebra is class~$\mathbf{L}$ ($r_{\max}=3$):
the Jacobi identity kills the quartic shadow on every primary line.
For a generic nontrivial DS reduction, the stress-tensor line is
required to carry a nonzero Virasoro quartic
$S_4 = 10/[c(5c{+}22)]$, away from the singular central charges
$c=0,-22/5$. The depth gap would then force
$r_{\max}=\infty$. Singular levels and degenerate central charges
require separate computation.

\begin{conjecture}[Shadow depth escalation under DS reduction]
\label{conj:ds-shadow-escalation-standalone}
Let $\fg$ be a simple Lie algebra, $f \in \fg$ a non-zero
nilpotent, and let $k$ be a non-critical level for which the
stress-tensor line of $\cW^k(\fg,f)$ is nondegenerate and has
$c\notin\{0,-22/5\}$. Then
\[
 r_{\max}\bigl(\cW^k(\fg, f)\bigr)
 \;\geq\;
 r_{\max}\bigl(V_k(\fg)\bigr)
 \;=\; 3.
\]
The conjecture asserts strictness for every non-zero nilpotent:
$r_{\max}(\cW^k(\fg, f)) = \infty$
\textup{(}class~$\mathbf{M}$\textup{)}.
\end{conjecture}

The conjecture asserts that DS reduction is a one-way door in
shadow complexity: once the nilpotent $f$ is turned on, the
tower remains infinite in the conjectural generic regime. The depth gap
(Proposition~\ref{prop:depth-gap-standalone}) reduces the
conjecture to a single computation: it suffices to show
$S_4 \neq 0$ on the stress-tensor primary line of
$\cW^k(\fg, f)$. The infinite tower then follows
automatically.

\begin{remark}[Evidence]\label{rem:ds-escalation-evidence}
Three independent lines of evidence support the conjecture.

\emph{(a)~Principal reductions.} For $f = f_{\mathrm{prin}}$
and the standard generic-level principal rows, the statement is proved: the principal
$\cW$-algebra $\cW_N = \cW^k(\mathfrak{sl}_N,
f_{\mathrm{prin}})$ contains the Virasoro algebra as a
subalgebra, and the Virasoro quartic $S_4 \neq 0$ on the
$T$-line is a universal computation independent of~$\fg$.

\emph{(b)~Non-principal reductions in type~$A$.} Existing
computations support the conjecture for several nilpotent orbits of
$\mathfrak{sl}_N$ at non-critical levels~$k$. The generic statement
remains conjectural outside the verified OPE tables.
The minimal nilpotent orbit of $\mathfrak{sl}_3$ (the
Bershadsky--Polyakov algebra $\cW^k(\mathfrak{sl}_3, f_{\min})$)
is the simplest non-principal example: it has a single
strong generator beyond the stress tensor, and $S_4 \neq 0$ is
verified by direct OPE computation.

\emph{(c)~Structural mechanism.} Every non-trivial
nilpotent~$f$ produces, via the BRST cohomology of the
associated quantum DS complex, at least one composite strong
generator whose self-OPE has pole order $\geq 4$. The pole
order of the composite is bounded below by the Sugawara
construction: the stress tensor $T$ always appears as a strong
generator, and $T(z)T(w)$ has a quartic pole. This quartic pole
creates $S_4 \neq 0$ on the stress-tensor primary line, and the
depth gap converts this single nonzero quartic into an infinite
tower.
\end{remark}


\subsection{The join-semilattice of shadow classes}
\label{ssec:shadow-semilattice}

The four shadow classes are invariants of individual algebras.
Their behaviour under algebraic operations requires a separate
argument. The split operation considered here is the direct sum of
cyclically admissible Lie conformal algebras with vanishing mixed
OPE, whose modular envelopes split as completed tensor products.

The answer is the maximum, and the proof is that the shadow
tower decomposes sector by sector when the cross-sector OPE is
regular. This gives the four shadow classes the structure of
a join-semilattice: a partially ordered set in which every pair
of elements has a least upper bound, realised concretely by
direct sum.

\begin{proposition}[Shadow classes form a join-semilattice]
\label{prop:shadow-semilattice-standalone}
The four shadow classes $\{\mathbf{G},\, \mathbf{L},\,
\mathbf{C},\, \mathbf{M}\}$ form a join-semilattice under the
total order $\mathbf{G} < \mathbf{L} < \mathbf{C} < \mathbf{M}$,
where the join is realised by direct sum of chiral algebras:
\begin{equation}\label{eq:shadow-join}
 \mathrm{class}(\cA_1 \oplus \cA_2)
 \;=\;
 \max\bigl(\mathrm{class}(\cA_1),\,
 \mathrm{class}(\cA_2)\bigr).
\end{equation}
Equivalently, the shadow depth satisfies
$r_{\max}(\cA_1 \oplus \cA_2)
= \max\bigl(r_{\max}(\cA_1),\, r_{\max}(\cA_2)\bigr)$.
\end{proposition}

\begin{proof}
Assume $\cA_1$ and $\cA_2$ are cyclically admissible Lie conformal
algebras with vanishing mixed OPE and split modular envelopes. The OPE
of $\cA_1 \oplus \cA_2$ decomposes into three sectors:
pure $\cA_1$, pure $\cA_2$, and mixed. The mixed sector has
regular OPE by hypothesis: fields from distinct factors are mutually
local with no singular terms.
Regular OPE contributes no shadow coefficients at any degree.
The shadow tower of $\cA_1 \oplus \cA_2$ is therefore the
direct sum of the towers of $\cA_1$ and $\cA_2$, and the shadow
depth is the maximum of the two.

The decisive case is $\mathbf{L} \oplus \mathbf{L}$: does
combining two affine Kac--Moody factors produce a quartic
shadow? The answer is no. The quartic obstruction $S_4$
is intrinsic to each factor's OPE. On the primary line of
$V_k(\fg_1)$, the Jacobi identity kills $S_4$; on the primary
line of $V_{k'}(\fg_2)$, the same. The mixed channel
$\fg_1 \otimes \fg_2$ contributes no singular OPE terms, so
no quartic shadow is born in the cross-sector. Therefore
$\mathbf{L} \oplus \mathbf{L} = \mathbf{L}$: the Jacobi identity
of each factor survives the direct sum.
\end{proof}

\begin{remark}[Monotonicity and its failure]
\label{rem:semilattice-limitations}
The join-semilattice property is specific to direct sums: the
functor $\oplus$ preserves the regularity of cross-sector OPE.
Constructions that create new singular terms between sectors
violate monotonicity and can raise the shadow class strictly.
Two examples illustrate the failure.

Drinfeld--Sokolov reduction sends class~$\mathbf{L}$ affine
input to class~$\mathbf{M}$ output in the verified principal rows
and in the generic regime predicted by
Conjecture~\ref{conj:ds-shadow-escalation-standalone}. The BRST
construction creates composite fields with new OPE singularities;
these singularities produce the quartic shadow that the Jacobi
identity had killed.

Coset construction: $\cA / \cA'$ can land in a higher class
than either factor. The Goddard--Kent--Olive coset of
$V_k(\mathfrak{sl}_2)$ by its Cartan subalgebra produces the
parafermion algebra, whose shadow depth depends on the coset
data rather than on the shadow classes of the two inputs.

Shadow depth is monotone under direct sum, anti-monotone under
no operation, and unconstrained under quotient or reduction.
The join-semilattice records the stable part of this structure;
the unstable part is the subject of
Conjecture~\ref{conj:ds-shadow-escalation-standalone}.
\end{remark}

\begin{thebibliography}{99}

\bibitem{Arakawa07}
T. Arakawa, \emph{Representation theory of W-algebras}, Invent.
Math. \textbf{169} (2007), no.~2, 219--320.

\bibitem{Flohr1996}
M. A. I. Flohr, \emph{On modular invariant partition functions of
conformal field theories with logarithmic operators}, Int. J. Mod.
Phys. A \textbf{11} (1996), 4147--4172.

\bibitem{FP03}
C. Faber and R. Pandharipande, \emph{Hodge integrals, partition
matrices, and the $\lambda_g$ conjecture}, Ann. of Math. \textbf{157}
(2003), no.~1, 97--124.

\bibitem{Frenkel-Kac-Wakimoto92}
I. B. Frenkel, V. G. Kac, and M. Wakimoto, \emph{Characters and
fusion rules for $W$-algebras via quantized Drinfeld--Sokolov
reduction}, Comm. Math. Phys. \textbf{147} (1992), no.~2, 295--328.

\bibitem{Gurarie1993}
V. Gurarie, \emph{Logarithmic operators in conformal field theory},
Nucl. Phys. B \textbf{410} (1993), 535--549.

\bibitem{KRW03}
V. G. Kac, S.-S. Roan, and M. Wakimoto, \emph{Quantum reduction for
affine superalgebras}, Comm. Math. Phys. \textbf{241} (2003),
307--342.

\end{thebibliography}

\ifdefined\shadowtowervtwoincluded
\else
\end{document}
\fi
